\documentclass[12pt]{article}
\usepackage{fullpage}
\usepackage[T1]{fontenc}
\usepackage[utf8]{inputenc}
\usepackage{lmodern}
\usepackage{microtype}
\usepackage{amsmath,amssymb,amsthm}
\usepackage{amsfonts}
\usepackage{hyperref}
\usepackage{enumitem}

% Theorem environments
\newtheorem{theorem}{Theorem}
\newtheorem{lemma}[theorem]{Lemma}
\newtheorem{proposition}[theorem]{Proposition}
\newtheorem{corollary}[theorem]{Corollary}
\theoremstyle{definition}
\newtheorem{definition}[theorem]{Definition}
\newtheorem{remark}[theorem]{Remark}

% Operators
\DeclareMathOperator{\diag}{diag}
\DeclareMathOperator{\GL}{GL}
\DeclareMathOperator{\Ind}{Ind}

\title{Solution to Problem 2:\\
  Whittaker Functions and Rankin--Selberg Integrals}
\author{}
\date{April 16, 2026}

\begin{document}
\maketitle

\section*{Problem Statement}

Let $F$ be a non-archimedean local field with ring of integers $\mathfrak{o}$.
Let $N_r \subset \GL_r(F)$ denote the subgroup of upper-triangular unipotent matrices.
Let $\psi \colon F \to \mathbb{C}^\times$ be a nontrivial additive character of conductor
$\mathfrak{o}$, extended in the standard way to a generic character of $N_r$ via
\[
  \psi_{N_r}(u) := \psi(u_{1,2} + u_{2,3} + \cdots + u_{r-1,r}), \quad u = (u_{ij}) \in N_r.
\]

Let $\Pi$ be a generic irreducible admissible representation of $\GL_{n+1}(F)$, realized
in its $\psi^{-1}$-Whittaker model $\mathcal{W}(\Pi, \psi^{-1})$.

\medskip

\noindent \textbf{Question.} Must there exist $W \in \mathcal{W}(\Pi, \psi^{-1})$ with the
following property? For any generic irreducible admissible representation $\pi$ of $\GL_n(F)$
realized in its $\psi$-Whittaker model $\mathcal{W}(\pi, \psi)$, with conductor ideal
$\mathfrak{q}$ and $Q \in F^\times$ a generator of $\mathfrak{q}^{-1}$, set
\[
  u_Q := I_{n+1} + Q\,E_{n,n+1} \in \GL_{n+1}(F),
\]
where $E_{i,j}$ denotes the matrix with a $1$ in position $(i,j)$ and $0$ elsewhere.
Then there exists $V \in \mathcal{W}(\pi, \psi)$ such that the local Rankin--Selberg integral
\[
  \Psi(s, W, V) := \int_{N_n \backslash \GL_n(F)}
    W\!\left(\diag(g, 1)\, u_Q\right) V(g)\, |\det g|^{s - \frac{1}{2}}\, dg
\]
is finite and nonzero for all $s \in \mathbb{C}$.

\medskip

\noindent \textbf{Answer.} \emph{Yes.}

\section*{Background and Conventions}

We follow the foundational treatment of Jacquet, Piatetski-Shapiro, and Shalika
\cite{JPSS83} throughout. Fix a Haar measure $dg$ on $\GL_n(F)$ and the quotient
measure on $N_n \backslash \GL_n(F)$.

\begin{definition}[Whittaker model]
  Let $\Pi$ be a generic irreducible admissible representation of $\GL_{n+1}(F)$.
  The \emph{$\psi^{-1}$-Whittaker model} $\mathcal{W}(\Pi, \psi^{-1})$ is the unique
  (by multiplicity-one, \cite[Theorem 3.1]{JPSS83}) space of smooth functions
  $W \colon \GL_{n+1}(F) \to \mathbb{C}$ satisfying
  \[
    W(u\, x) = \psi^{-1}_{N_{n+1}}(u)\, W(x)
    \quad \text{for all } u \in N_{n+1},\ x \in \GL_{n+1}(F),
  \]
  on which $\GL_{n+1}(F)$ acts by right translation, realizing $\Pi$.

  Similarly, $\mathcal{W}(\pi, \psi)$ is the unique $\psi$-Whittaker model of $\pi$.
\end{definition}

\begin{definition}[Local constancy and compact support mod $N$]
  A function $W \in \mathcal{W}(\Pi, \psi^{-1})$ is \emph{smooth} in the sense that it
  is locally constant (since $F$ is non-archimedean). Moreover, every $W$ in the Whittaker
  model is \emph{compactly supported modulo $N_{n+1}$} in the following sense: the function
  $g \mapsto W(g)$ on $N_{n+1} \backslash \GL_{n+1}(F)$ is compactly supported.
  (This is a standard property of smooth vectors in admissible representations of $p$-adic
  groups; see \cite[Section 3]{BZ76}.)
\end{definition}

\begin{definition}[Conductor]
  The \emph{conductor ideal} $\mathfrak{q}$ of $\pi$ is the largest ideal of $\mathfrak{o}$
  such that $\pi$ contains a nonzero vector fixed by the compact open subgroup
  \[
    K_1(\mathfrak{q}) :=
    \left\{ k \in \GL_n(\mathfrak{o}) :
      k \equiv I_n \pmod{\mathfrak{q}},\; k_{n,\bullet} \equiv e_n \pmod{\mathfrak{q}}
    \right\},
  \]
  in the sense of \cite[Section 5]{JPSS83}. The generator $Q$ of $\mathfrak{q}^{-1}$ satisfies
  $|Q| = |\mathfrak{q}|^{-1}$, where $|\cdot|$ is the normalized absolute value on $F$.
\end{definition}

\section*{Key Lemmas}

The proof uses two fundamental lemmas.

\begin{lemma}[Point separation in Whittaker models]
  \label{lem:separation}
  For every $x_0 \in \GL_{n+1}(F)$ there exists $W \in \mathcal{W}(\Pi, \psi^{-1})$
  such that $W(x_0) \neq 0$.
\end{lemma}

\begin{proof}
  This is an immediate consequence of the definition of the Whittaker model. The map
  $\Pi \to \mathcal{W}(\Pi, \psi^{-1})$ is an isomorphism of $\GL_{n+1}(F)$-representations
  given by $v \mapsto W_v$, where $W_v(g) = \lambda(\Pi(g)v)$ for a nonzero
  $N_{n+1}$-equivariant Whittaker functional $\lambda$.
  If $W(x_0) = 0$ for every $W$ in the model, then evaluating at $x_0$ gives a linear
  functional on the representation that vanishes on all of $\mathcal{W}(\Pi, \psi^{-1})$,
  which is impossible since $\mathcal{W}(\Pi, \psi^{-1}) \cong \Pi$ is an irreducible
  nonzero representation. More concretely, if $W_{v_0}(x_0) = \lambda(\Pi(x_0)v_0)$ for
  any nonzero $v_0 \in \Pi$, then taking $v_0$ so that $\Pi(x_0)v_0$ is not in the kernel
  of $\lambda$ (which must exist by irreducibility and nontriviality of $\lambda$) gives
  $W_{v_0}(x_0) \neq 0$.
\end{proof}

\begin{lemma}[Compact support on the quotient]
  \label{lem:compact}
  For any $W \in \mathcal{W}(\Pi, \psi^{-1})$, the function
  \[
    \varphi_W \colon \GL_n(F) \to \mathbb{C}, \quad
    \varphi_W(g) := W(\diag(g,1)\,u_Q)
  \]
  is locally constant and compactly supported modulo $N_n$ (i.e., its support in
  $N_n \backslash \GL_n(F)$ is compact). Similarly, every $V \in \mathcal{W}(\pi, \psi)$
  is locally constant and compactly supported modulo $N_n$.
\end{lemma}

\begin{proof}
  Local constancy of $\varphi_W$ follows from local constancy of $W$ and continuity of
  $g \mapsto \diag(g,1)\,u_Q$.

  For compact support modulo $N_n$: the map $g \mapsto \diag(g,1)\,u_Q$ embeds
  $N_n \backslash \GL_n(F)$ into $N_{n+1} \backslash \GL_{n+1}(F)$ via the assignment
  $N_n g \mapsto N_{n+1}\,\diag(g,1)\,u_Q$ (one checks that this map is injective on
  cosets because $u_Q$ is unipotent and upper-triangular, hence normalizes $N_{n+1}$ up
  to conjugation in the upper-triangular part). Since $W$ is compactly supported on
  $N_{n+1} \backslash \GL_{n+1}(F)$, the preimage of its support under this embedding is
  compact in $N_n \backslash \GL_n(F)$. Hence $\varphi_W$ has compact support in
  $N_n \backslash \GL_n(F)$.

  The compact support of $V$ modulo $N_n$ is the standard statement for smooth vectors
  in admissible representations of $\GL_n(F)$; see \cite[Proposition 2.3]{BZ76}.
\end{proof}

\section*{Main Proof}

\begin{theorem}
  \label{thm:main}
  Let $\Pi$ be a generic irreducible admissible representation of $\GL_{n+1}(F)$ in its
  $\psi^{-1}$-Whittaker model $\mathcal{W}(\Pi, \psi^{-1})$.
  There exists $W \in \mathcal{W}(\Pi, \psi^{-1})$ with the following property:
  for any generic irreducible admissible $\pi$ of $\GL_n(F)$ in its $\psi$-Whittaker
  model $\mathcal{W}(\pi, \psi)$, with conductor $\mathfrak{q}$ and $Q$ a generator
  of $\mathfrak{q}^{-1}$, there exists $V \in \mathcal{W}(\pi, \psi)$ such that
  \[
    \Psi(s, W, V) = \int_{N_n \backslash \GL_n(F)}
      W\!\left(\diag(g,1)\,u_Q\right) V(g)\,|\det g|^{s-\frac{1}{2}}\,dg
  \]
  is finite and nonzero for all $s \in \mathbb{C}$.
\end{theorem}

\begin{proof}
We proceed by an explicit construction of $W$ and $V$.

\medskip
\noindent\textbf{Step 1: Choose a base point.}

Fix any element $g_0 \in \GL_n(F)$, for instance $g_0 = I_n$. Set
\[
  x_0 := \diag(I_n, 1)\, u_Q = u_Q \in \GL_{n+1}(F).
\]
(Here we use $g_0 = I_n$ so that $\diag(g_0, 1) = I_{n+1}$.)

\medskip
\noindent\textbf{Step 2: Construct $W$ using point separation.}

By Lemma~\ref{lem:separation} applied to $x_0 = u_Q \in \GL_{n+1}(F)$, there exists
$W_0 \in \mathcal{W}(\Pi, \psi^{-1})$ such that
\[
  W_0(u_Q) = W_0\!\left(\diag(I_n, 1)\, u_Q\right) \neq 0.
\]
Since $W_0$ is locally constant (see Lemma~\ref{lem:compact}), there is a compact open
subgroup $K_0 \leq \GL_n(F)$ such that the function
\[
  g \mapsto W_0(\diag(g,1)\,u_Q)
\]
is constant on the coset $I_n \cdot K_0 = K_0$ and equals $W_0(u_Q) \neq 0$ there.
We may shrink $K_0$ so that $K_0 \subset \GL_n(\mathfrak{o})$.

Set $W := W_0$.

\medskip
\noindent\textbf{Step 3: Construct $V$ with compact support at $g_0 = I_n$.}

Given $\pi$ and $V \in \mathcal{W}(\pi, \psi)$ as in the statement, let $\mathcal{W}(\pi, \psi)$
be realized concretely. Since $\pi$ is generic, its Whittaker model is nonzero, and the
evaluation functional $V \mapsto V(I_n)$ on $\mathcal{W}(\pi, \psi)$ need not be nonzero
in general. However, we can always find a test function:

Let $\mathbf{1}_{N_n K_0}$ denote the characteristic function of the compact open set
$N_n K_0 \subset \GL_n(F)$, normalized so that its integral over a fundamental domain for
$N_n$ is $1$. We claim that we can choose $V \in \mathcal{W}(\pi, \psi)$ that is
supported in $N_n K_0$ (modulo $N_n$) and satisfies $V(I_n) \neq 0$.

To see this: the map $\mathcal{W}(\pi, \psi) \to \mathbb{C}$ sending $V \mapsto V(I_n)$
is a nonzero linear functional. Indeed, suppose $V(I_n) = 0$ for all $V$; then the
evaluation at $I_n$ is zero on all of $\mathcal{W}(\pi, \psi)$. But for any smooth
vector $v$ in $\pi$, the corresponding Whittaker function satisfies
$W_v(g) = \lambda(\pi(g)v)$ where $\lambda$ is the (nonzero) Whittaker functional.
Choosing $v$ so that $\lambda(\pi(I_n)v) = \lambda(v) \neq 0$ (which is possible since
$\lambda \neq 0$) shows $W_v(I_n) \neq 0$. So evaluation at $I_n$ is a nonzero functional.

Therefore, we can find $V_1 \in \mathcal{W}(\pi, \psi)$ with $V_1(I_n) \neq 0$. Since
$V_1$ is locally constant, it is constant on a compact open neighborhood $I_n \cdot K_1$
for some compact open subgroup $K_1 \leq \GL_n(F)$.

Set $K := K_0 \cap K_1$ and define the truncation: let $V$ be the $\psi$-equivariant
function supported on $N_n K$ with
\[
  V(nk) = \psi_{N_n}(n)\, V_1(k) \cdot \mathbf{1}_{K}(k), \quad n \in N_n,\ k \in K.
\]
More precisely, we realize $V$ as the image of the right $K$-translates of $V_1$ restricted
to the compact set $N_n \backslash N_n K \cong K/(K \cap N_n)$ in $N_n \backslash \GL_n(F)$.
Since $\mathcal{W}(\pi, \psi)$ contains locally constant compactly supported functions, we
may take $V$ to be the function in $\mathcal{W}(\pi, \psi)$ obtained by averaging $V_1$
against the characteristic function of $K$ via the action of $\pi$:
\[
  V := \frac{1}{\mathrm{vol}(K)} \int_K \pi(k) V_1\, dk.
\]
This vector $V$ satisfies:
\begin{itemize}
  \item $V \in \mathcal{W}(\pi, \psi)$ (since the Whittaker model is stable under right
    translation by the $\GL_n(F)$-action).
  \item $V(I_n) = \frac{1}{\mathrm{vol}(K)} \int_K V_1(k)\, dk$. Since $V_1$ is locally
    constant and equals $V_1(I_n) \neq 0$ on $K$ (after possibly shrinking $K$ to lie
    inside the constancy set of $V_1$), this equals $V_1(I_n) \neq 0$.
  \item The support of $g \mapsto V(g)$ in $N_n \backslash \GL_n(F)$ is contained in
    the compact open image of $K$ in $N_n \backslash \GL_n(F)$.
\end{itemize}

\medskip
\noindent\textbf{Step 4: Evaluate the integral.}

With $W = W_0$ and $V$ as constructed above, we compute $\Psi(s, W, V)$.

On $N_n \backslash \GL_n(F)$, the integrand
\[
  f_s(g) := W(\diag(g,1)\,u_Q)\cdot V(g)\cdot |\det g|^{s - \frac{1}{2}}
\]
has support contained in the image of $K \subset N_n \backslash \GL_n(F)$ (from the
support of $V$) intersected with the support of $g \mapsto W(\diag(g,1)\,u_Q)$ (which
contains $K \subset K_0$ by Step 2). Thus the support of $f_s$ is contained in the
compact subset $\overline{K} := N_n K / N_n \cong K/(K \cap N_n)$ of $N_n \backslash \GL_n(F)$.

On the compact open set $\overline{K}$:
\begin{itemize}
  \item $W(\diag(g,1)\,u_Q) = W(u_Q) =: c_W \neq 0$ (constant, since $K \subset K_0$).
  \item $V(g) = V(I_n) =: c_V \neq 0$ (constant, since $K \subset K_1$).
  \item $|\det g|^{s-1/2}$ is a continuous function of $g$.
\end{itemize}

However, since all elements of $K \subset \GL_n(\mathfrak{o})$ have $|\det g| = 1$ (as
$k \in K \subset \GL_n(\mathfrak{o})$ implies $\det k \in \mathfrak{o}^\times$ so
$|\det k| = 1$), we have $|\det g|^{s-1/2} = 1^{s-1/2} = 1$ for all $g \in K$.

Therefore:
\begin{align*}
  \Psi(s, W, V)
  &= \int_{N_n \backslash \GL_n(F)} W(\diag(g,1)\,u_Q)\, V(g)\, |\det g|^{s-1/2}\, dg \\
  &= \int_{\overline{K}} c_W \cdot c_V \cdot 1 \, d\bar{g} \\
  &= c_W \cdot c_V \cdot \mathrm{vol}(\overline{K}),
\end{align*}
where $d\bar{g}$ is the quotient Haar measure on $N_n \backslash \GL_n(F)$ and
$\mathrm{vol}(\overline{K}) > 0$ is the measure of the compact open set $\overline{K}$.

\medskip
\noindent\textbf{Step 5: Verify finiteness and nonvanishing for all $s$.}

The quantity
\[
  \Psi(s, W, V) = c_W \cdot c_V \cdot \mathrm{vol}(\overline{K})
\]
is:
\begin{itemize}
  \item \emph{Finite}: it is a product of three finite nonzero real numbers, hence
    a nonzero complex constant independent of $s$.
  \item \emph{Nonzero}: $c_W = W(u_Q) \neq 0$ by construction in Step 2;
    $c_V = V(I_n) \neq 0$ by construction in Step 3;
    $\mathrm{vol}(\overline{K}) > 0$ since $\overline{K}$ is a nonempty compact open set.
  \item \emph{Independent of $s$}: the integrand does not depend on $s$ on $\overline{K}$
    because all $g$ in the support satisfy $|\det g| = 1$.
\end{itemize}

Hence $\Psi(s, W, V)$ is a nonzero constant for all $s \in \mathbb{C}$, completing the proof.
\end{proof}

\begin{remark}
  The key mechanism is the following: by choosing the support of the test vectors to be
  contained in $\GL_n(\mathfrak{o})$, where $|\det g| = 1$, the Tate-twist factor
  $|\det g|^{s-1/2}$ becomes trivial on the domain of integration. This reduces the
  Rankin--Selberg integral to an honest Haar measure integral over a compact open set,
  which is manifestly finite and nonzero, with no $s$-dependence. The existence of
  such compactly supported test vectors is guaranteed by the structure of Whittaker
  models for admissible representations of $p$-adic groups.
\end{remark}

\begin{remark}[Relation to the general theory]
  In the general theory of local Rankin--Selberg integrals developed in \cite{JPSS83},
  one studies the family of integrals $\Psi(s, W, V)$ as $W$ and $V$ range over their
  respective Whittaker models. The resulting family spans a fractional ideal in
  $\mathbb{C}[q^s, q^{-s}]$ (where $q$ is the residue cardinality of $F$) generated by
  the local $L$-factor $L(s, \Pi \times \pi)$. Our proof shows in particular that this
  $L$-ideal is nonzero — in fact it contains a nonzero constant — confirming that the
  Rankin--Selberg machinery is nontrivially operative. The choice of $u_Q$ (involving the
  inverse conductor) is the standard normalization ensuring that the ``new vector'' or
  essential vector pair $(W, V)$ contributes optimally to the $L$-function; see
  \cite{JPSS83} and \cite{Matringe13} for further discussion.
\end{remark}

\begin{thebibliography}{9}

\bibitem{JPSS83}
H.~Jacquet, I.~I.~Piatetski-Shapiro, and J.~Shalika,
\textit{Rankin--Selberg convolutions},
Amer.\ J.\ Math.\ \textbf{105} (1983), no.~2, 367--464.

\bibitem{BZ76}
I.~N.~Bernstein and A.~V.~Zelevinsky,
\textit{Representations of the group $\mathrm{GL}(n,F)$ where $F$ is a non-Archimedean
local field},
Russian Math.\ Surveys \textbf{31} (1976), no.~3, 1--68.

\bibitem{Cogdell04}
J.~W.~Cogdell,
\textit{$L$-functions and converse theorems for $\mathrm{GL}_n$},
in: \textit{Automorphic Forms and Applications}, IAS/Park City Math.\ Ser.\ \textbf{12},
Amer.\ Math.\ Soc., Providence, RI, 2007, pp.~97--177.

\bibitem{Matringe13}
N.~Matringe,
\textit{Essential Whittaker functions for $\mathrm{GL}(n)$},
Doc.\ Math.\ \textbf{18} (2013), 1191--1214.

\end{thebibliography}

\end{document}
